% SPDX-FileCopyrightText: Copyright (C) Nile Jocson <atraphaxia@gmail.com>
% SPDX-License-Identifier: MPL-2.0



\section{Notes}

\anotegroup{Quadratics}{
	\anote{Definition}{qd}{%
		Let $a, b, c \in \mathbb{R}$ and $a \neq 0$. A quadratic equation in standard form is:
		\[ax^2 + bx + c = 0\]
		These equations have either one or two solutions. Note that the square of a number cannot be negative.
	}
	\anote{Discriminant}{qdi}{%
		Given a quadratic in standard form, its discriminant is:
		\[b^2 - 4ac\]
		\begin{enumerate}
			\item If this is greater than zero, the quadratic has two real solutions.
			\item If this is equal to zero, the quadratic has one real solution.
			\item If this is less than zero, the quadratic has two complex solutions which are conjugates of each other.
		\end{enumerate}
	}
	\anote{Quadratic Formula}{qqf}{%
		Given a quadratic in standard form, the formula to solve for its solutions is:
		\[\dfrac{-b \pm \sqrt{b^2 - 4ac}}{2a}\]
	}
	\anote{Factoring}{qf}{%
		Given a quadratic in standard form, you may factor it in order to solve for its solutions. For example:
		\begin{align*}
			4x^2 + 4x + 1 &= 0 \\
			{(2x + 1)}^2 &= 0 \\
			\Aboxed{x &= -\dfrac{1}{2}}
		\end{align*}
	}
	\anote{Equations in Quadratic Form}{qeiqf}{%
		Given the equation:
		\[2{\ab(x + \dfrac{1}{x})}^2 + \ab(x + \dfrac{1}{x}) - 10 = 0\]
		While not a quadratic equation, we may rewrite it as such:
		\[2t^2 + t - 10 = 0 \text{, where } t = x + \dfrac{1}{x}\]
		Solve for the values of $t$ as usual, then solve for $x$ using $t$. For example, using the equation above:
		\begin{align*}
			2{\ab(x + \dfrac{1}{x})}^2 + \ab(x + \dfrac{1}{x}) - 10 &= 0
			\intertext{$t = x + \dfrac{1}{x}$}
			2t^2 + t - 10 &= 0 \\
			2t^2 + 5t - 4t - 10 &= 0 \\
			(t - 2)(2t + 5) &= 0 \\
			t &\in \ab\{-\dfrac{5}{2}, 2\}
			\intertext{Solve for $x$, $t = -\dfrac{5}{2}$.}
			x + \dfrac{1}{x} &= -\dfrac{5}{2}
			\intertext{$x \neq 0$}
			2x^2 + 5x + 2 &= 0 \\
			2x^2 + 4x + x + 2 &= 0 \\
			(2x + 1)(x + 2) &= 0 \\
			x &\in \ab\{-2, -\dfrac{1}{2}\}
			\intertext{Solve for $x$, $t = 2$.}
			x + \dfrac{1}{x} &= 2
			\intertext{$x \neq 0$}
			x^2 - 2x + 1 &= 0 \\
			{(x - 1)}^2 &= 0 \\
			x &= 1 \\
			\Aboxed{x &\in \ab\{-2, -\dfrac{1}{2}, 1\}}
		\end{align*}
	}
}

\anotegroup{Radicals}{
	\anote{Negative Numbers}{rnn}{%
		Note that $\sqrt[n]{x}$ cannot be negative for even values of $n$.
	}
	\anote{Extraneous Solutions}{res}{%
		It is required to check each solution when exponentiating both sides of an equation, as these may
		be extraneous solutions, which are not solutions to the original equation.
	}
}

\anotegroup{Absolute Values}{
	\anote{Definition}{avd}{
		\[|x| = \begin{cases}
			\phantom{-}x,& x \geq 0 \\
			-x,& x < 0
		\end{cases}\]
		Note that the absolute value of a number cannot be negative.
	}
	\anote{Of Products}{avop}{
		\[|xy| = |x||y|\]
	}
	\anote{Of Quotients}{avoq}{
		\[\ab|\dfrac{x}{y}| = \dfrac{|x|}{|y|}\]
	}
}

\anotegroup{Inequalities}{
	\anote{Quadratics}{iq}{%
		In order to solve quadratic inequalities, first rewrite it in standard form, factor, identify critical numbers
		and undefined points, then create a table of signs.
	}
	\anote{Absolute Values}{iav}{%
		In order to solve inequalities with absolute values, take note of the following:
		\begin{enumerate}
			\item $|x| < a$ if and only if $-a < x < a$.
			\item $|x| > a$ if and only if $x < -a$ or $x > a$.
		\end{enumerate}
	}
}
